Late-Time Geometric Torque as Dynamical Expression of the Nested Phase Structure

The elevation of the present-day Hubble constant from the bosonic floor of \(70\,\mathrm{km\,s^{-1}Mpc^{-1}}\) into the observationally preferred window of \(73\)–\(74\,\mathrm{km\,s^{-1}Mpc^{-1}}\) is not an independent phenomenological adjustment. It is the temporal manifestation of a single, spatially organized phase structure that already exists, fully formed, inside the mid-shell and outer shells of the \(360^\circ\) Sphere Radians Representation. The late-time torque \(\Delta H_{\rm torque}\) is therefore the dynamical reading of the same geometric architecture that appears statically as the Quaternionic Orders \(1\)–\(12\) and the Octonionic Orders \(10/15/30\).

At the algebraic root stands the ordered pair \(i:=(0,1)\). This choice fixes the counterclockwise-positive orientation of every subsequent winding. Continuous accumulation of that orientation generates the maximal spiral of winding number \(w=3.25\). Discrete angular sampling of the same spiral isolates four privileged phases—\(\psi\), \(\theta_{\rm eff}\), \(\Delta\phi_{\rm torque}\) and \(\pm\arctan(2\pi)\)—that satisfy the exact identity
\[
\sin(\Delta\phi_{\rm torque})=\cos\psi\approx0.988721.
\]
These four phases are the generators of the Cosmological Clock’s torque term. When the suppression functions and scale-factor weighting are applied, they produce a positive contribution that lifts the bosonic baseline by approximately \(3.2\)–\(3.7\,\mathrm{km\,s^{-1}Mpc^{-1}}\) at \(z=0\), yielding a total \(H_0\) in the \(73\)–\(74\) range.

The identical phases organize the spatial geometry of the nested shells. In the Quaternionic Baseline the six-fold radial symmetry of the starred polyhedra is the discrete embodiment of the continuous winding whose characteristic increments are precisely those four angles. Orders \(1\) through \(4\) construct the hexagram lattice step by step; Order \(12\) densifies the same lattice into a multi-layered mid-shell whose interlaced edges encode the non-commutativity of \(\mathbb{H}\). Every edge and every vertex continues to point along directions fixed by \(\psi\), \(\theta_{\rm eff}\), \(\Delta\phi_{\rm torque}\) and the terminal ray. The central real-scalar node that survives in every order is the geometric counterpart of the bosonic backbone that remains after the torque has been isolated.

The outer shells extend the same logic into the octonionic regime. Orders \(10\), \(15\) and \(30\) populate successive concentric Petrie rings whose angular sampling approximates the \(E_8\) root system. Because the rings are built from the identical angular set, the phase slip responsible for the late-time torque is already latent in the highest shell of the vacuum geometry. The residual Gaussian compression at the Transition Lock (\(t=4/7\)) further modulates the amplitude of that phase slip, providing the redshift-dependent suppression that keeps the torque negligible at early times and dominant only in the late universe.

Thus the numerical value of \(\Delta H_{\rm torque}\) is not an extra degree of freedom inserted by hand. It is the projection, onto the expansion history, of a phase structure that has already been fixed by the requirement of consistency across the Quaternionic mid-shell and the Octonionic outer shells. The same geometric object appears once as a static, ordered polyhedral lattice on the unit sphere and once as a dynamical contribution to \(H(z)\). The elevation of \(H_0\) from \(70\) to the \(73\)–\(74\) window is simply the late-time reading of that object.

This identification supplies a stringent coherence condition. Any modification of the angular parameters that alters the torque must simultaneously deform the nodal structure of the mid-shell starred polyhedra and the angular density of the outer \(E_8\) rings. Conversely, any geometric refinement of the nested shells that fails to recover the original four phases cannot reproduce the observed late-time enhancement of \(H_0\). The dynamical and the geometric descriptions are therefore locked together: the torque is the mid-shell and the outer shells written in the language of cosmic expansion.

In this sense the entire hierarchy—from the ordered pair \(i:=(0,1)\) through the progressive Quaternionic and Octonionic baselines to the present-day value of the Hubble constant—constitutes a single closed narrative. The late-time torque is not an addition to the geometry; it is the geometry itself, expressed as expansion.
